\chapter{Precision, calibration, and experimental design}
\label{precision:chapter}

Exact identifiability asks whether some accessible record law changes when the
scientific target changes. Precision asks how much evidence is needed to resolve
that change. These questions have different answers. A nonzero rational contrast
can be arbitrarily small, an invertible response matrix can be arbitrarily poorly
conditioned, and a short distinguishing experiment can require many repetitions.
This chapter supplies upper and lower guarantees that connect the exact access
layer to decisions a scientist can use.

Throughout, $\theta$ denotes a fixed physical model, $g(\theta)$ the target, and
$P_\theta^\pi$ the law of the complete accessible record under a common permitted
protocol $\pi$. A resource contract specifies preparation, actions, stopping,
and the meaning of every budget. The general statistical ingredients below are
established: binary testing, optimal recovery, concentration, and belief-state
dynamic programming. Their role here is to make the interfaces explicit.

\section{What the best possible precision means}

For probability laws $P,Q$ on a common record space, use the convention
\[
 \operatorname{TV}(P,Q)=\sup_A|P(A)-Q(A)|
 =\tfrac12\int|p-q|\,d\mu,
\]
where $\mu$ dominates both laws. If the two models are known and equally likely,
the smallest average probability of choosing the wrong model is
\begin{equation}
 r^*(P,Q)=\frac{1-\operatorname{TV}(P,Q)}2.
 \label{precision:tv-risk}
\end{equation}
Indeed, choose the first model on $\{p\geq q\}$; the resulting error is
$\frac12\int\min(p,q)\,d\mu$, and pointwise minimization shows that no other
decision rule has smaller error. This is an average error for specified priors.
It is not automatically a bound on each of the two conditional errors.

Let $\Pi_C$ be the declared class of protocols admissible at resource budget $C$.
With the access setting fixed, abbreviate the frontier of
Equation~\eqref{foundations:frontier} by
\begin{equation}
 D_C(\theta_0,\theta_1)=
 \sup_{\pi\in\Pi_C}\operatorname{TV}(P_{\theta_0}^\pi,P_{\theta_1}^\pi).
 \label{precision:resource-discrimination}
\end{equation}
For known equal-prior alternatives, the infimum achievable average error is
$(1-D_C)/2$. A supremum need not be attained by a protocol. If the budget classes
are nested, $D_C$ is nondecreasing. No shape, universal energy scale, or efficient
algorithm for this function follows from its definition.

\begin{proposition}[A resource-dependent lower bound for real targets]
\label{precision:two-point}
Let $g$ be real valued, and let an estimator $T$ use the record of any
$\pi\in\Pi_C$. The policy, its randomization rule, and the estimator must be the
same under both alternatives. Then
\begin{equation}
 \max_{j=0,1}\mathbb E_{\theta_j}^{\pi}|T-g(\theta_j)|
 \geq \frac{|g(\theta_0)-g(\theta_1)|}{2}
 \bigl[1-\operatorname{TV}(P_{\theta_0}^\pi,P_{\theta_1}^\pi)\bigr].
 \label{precision:two-point-fixed}
\end{equation}
Consequently, for any parameter family $\Theta$,
\begin{equation}
 \inf_{\pi\in\Pi_C,\,T}\sup_{\theta\in\Theta}
 \mathbb E_\theta^\pi|T-g(\theta)|
 \geq \sup_{\theta_0,\theta_1\in\Theta}
 \frac{|g(\theta_0)-g(\theta_1)|}{2}
 [1-D_C(\theta_0,\theta_1)].
 \label{precision:two-point-resource}
\end{equation}
Expectations may be infinite. The bound is not asserted to be attainable.
\end{proposition}
\begin{proof}
Write $g_j=g(\theta_j)$ and $p_j=dP_{\theta_j}^\pi/d\mu$. The triangle inequality
gives, pointwise,
\[
 p_0|T-g_0|+p_1|T-g_1|
 \geq\min(p_0,p_1)(|T-g_0|+|T-g_1|)
 \geq |g_0-g_1|\min(p_0,p_1).
\]
Integration and $\int\min(p_0,p_1)=1-\operatorname{TV}(P_0,P_1)$ prove
\eqref{precision:two-point-fixed}. Independent estimator randomization is a
common probability kernel and the same argument applies after adjoining its
random seed. For every fixed pair, replacing the policy's total variation by
its supremum over $\Pi_C$ weakens the bound uniformly in $\pi$. Taking the
infimum on the left and then the supremum over pairs proves
\eqref{precision:two-point-resource}; no minimax interchange is used.
\end{proof}

This is the standard two-point method associated with Le Cam
\cite{caLeCamYang2000}. It identifies an important failure case: an invisible
target difference has $D_C=0$ at every allowed budget, so more computation on
the same records cannot remove its ambiguity. If invisible pairs have arbitrarily
large finite target separation, the worst absolute-error risk is infinite.
The conclusion follows by taking those finite pairs, without undefined products
involving infinity. A lower bound caused by restricted access is different from
a claim that all information about the system has a universal cap.

\section{Linear intervention responses: from rank to stable recovery}

Suppose $\theta\in\Theta\subset\mathbb R^p$ and a known matrix
$F\in\mathbb R^{m\times p}$ represents intervention responses. Two physical
assumptions are needed before its algebra describes the apparatus:
\begin{enumerate}
 \item If $F\theta=F\theta'$, every permitted complete record law agrees.
 \item If every permitted complete record law agrees, then $F\theta=F\theta'$.
\end{enumerate}
The first is law factorization; the second is access sufficiency. A hidden
generator's dependence on $F\theta$ may establish the first while leaving the
second unproved for a coarse detector. Unknown preparation or a
parameter-dependent detector can invalidate either identification argument.

For example, controlled continuous-time generators with known entries
$q_0+F\theta$ give factorization when preparations and readout are known and
parameter independent. Access sufficiency follows in an idealized setting where
every state is preparable and resolved, and transition laws are accessible at
arbitrarily short times: the derivative at zero identifies each generator entry.
Finite bandwidth and finite data require additional bounds; that idealization
does not provide them.

\begin{proposition}[A target-specific row-space criterion]
\label{precision:row-space}
Assume both conditions above and a nonempty open domain $\Theta$. The linear
target $g(\theta)=b^{\mathsf T}\theta$ is identifiable if and only if
$b\in\operatorname{row}(F)$, equivalently $F^{\mathsf T}w=b$ for some $w$.
\end{proposition}
\begin{proof}
Such a $w$ yields $g(\theta)=w^{\mathsf T}F\theta$, so access sufficiency proves
the claim. Conversely, if $b$ is outside the row space, there is
$h\in\ker(F)$ with $b^{\mathsf T}h\ne0$. At an interior point $\theta_*$,
both $\theta_*+th$ and $\theta_*-th$ are allowed for sufficiently small $t>0$.
Factorization makes their laws identical, while their targets differ.
\end{proof}

On a relatively open affine parameter domain $\theta_*+V$, the correct condition
is that $b$ annihilate $V\cap\ker F$. Known constraints can therefore resolve
ambiguities that remain on the full open domain. Full parameter reconstruction
requires more than reconstruction of a particular scientific target.

Now suppose the measured or statistically estimated response satisfies
\begin{equation}
 \widehat y=F\theta+e,\qquad |e_i|\leq\epsilon_i,
 \label{precision:response-error}
\end{equation}
where every $\epsilon_i$ is known and nonnegative. Any $F^{\mathsf T}w=b$ gives
\begin{equation}
 |w^{\mathsf T}\widehat y-b^{\mathsf T}\theta|
 \leq \sum_i\epsilon_i|w_i|.
 \label{precision:linear-upper}
\end{equation}
No independence of coordinate errors is needed. If the entire box is certified
with probability at least $1-\delta$, this deterministic conclusion has at least
that coverage. How the box is obtained remains a separate experimental question.

\begin{theorem}[Global deterministic minimax recovery]
\label{precision:linear-minimax}
Assume the parameter domain is the entire positive orthant
$\mathbb R^p_{>0}$, $F$ is known, and every error vector in the full Cartesian
box $[-\epsilon,\epsilon]$ is admissible. If $b\in\operatorname{row}(F)$, then
\begin{align}
 &\inf_T\sup_{\theta>0,\,|e|\leq\epsilon}
       |T(F\theta+e)-b^{\mathsf T}\theta| \nonumber\\
 &\hspace{8mm}=\min_{F^{\mathsf T}w=b}\sum_i\epsilon_i|w_i|
   =\max_{|Fh|\leq\epsilon}b^{\mathsf T}h.
 \label{precision:linear-duality}
\end{align}
The extrema on the right exist, including when some error radii vanish.
If $b\notin\operatorname{row}(F)$, the recovery risk is infinite, the middle
problem is infeasible, and the dual supremum is unbounded.
\end{theorem}
\begin{proof}
Equation~\eqref{precision:linear-upper} proves the upper bound. For any dual-feasible
$h$, choose a componentwise center $\theta_*>|h|$. The two positive parameters
$\theta_*\pm h$, with errors $\mp Fh$, produce the same observation
$\widehat y=F\theta_*$. Their target separation is $2|b^{\mathsf T}h|$.
Every estimate at that common observation errs by at least $|b^{\mathsf T}h|$
on one of them. The feasible set for $h$ is symmetric, giving the displayed
dual lower bound.

Write $w=u-v$, $u,v\geq0$. The middle value equals the linear program
\[
 \min_{u,v\geq0}\epsilon^{\mathsf T}(u+v)
 \quad\hbox{subject to }F^{\mathsf T}u-F^{\mathsf T}v=b.
\]
Its dual is exactly $\max b^{\mathsf T}h$ subject to
$Fh\leq\epsilon$ and $-Fh\leq\epsilon$. The primal is feasible with a finite
nonnegative infimum, and the dual is feasible at zero. Finite-dimensional linear
programming duality supplies attainment and equality \cite{caBoyd2004}.
If $b$ is outside the row space, scale an $h\in\ker F$ with
$b^{\mathsf T}h>0$ and repeat the common-observation construction to obtain
arbitrarily large error lower bounds.
\end{proof}

An independently checkable certificate consists of $w,h$ satisfying
$F^{\mathsf T}w=b$, $|Fh|\leq\epsilon$, and
$\sum_i\epsilon_i|w_i|=b^{\mathsf T}h$. The checker need not trust the optimizer.
This is a structured, globally worst-case specialization of the compatible-set
geometry in Proposition~\ref{foundations:radius}. A particular observed
$\widehat y$ may give a smaller feasible target interval. Bounded parameter
domains or correlated error constraints can also lower the minimax value;
their replacement by the full positive orthant or full error box changes the
problem being solved.

\subsection{A calibrated intervention example}

Two unresolved positive channel rates have baseline response
$y_0=\theta_1+\theta_2$. A calibrated intervention produces
$y_1=r\theta_1+\theta_2$, with known $r\geq0$. Thus
\[
 F_r=\begin{pmatrix}1&1\\r&1\end{pmatrix},\qquad
 \theta_1=\frac{y_1-y_0}{r-1}\quad(r\ne1).
\]
At equal response-error radius $\epsilon$, the exact global minimax error for
$\theta_1$ is $2\epsilon/|r-1|$. For the primal choose
$w=(-1,1)/(r-1)$; a matching dual is
$h=\epsilon(2,-r-1)/|r-1|$. Its two response coordinates have magnitudes
$\epsilon$, and its first component equals the claimed error.
At $\epsilon=0.01$, responses $r=1/2,0.9,0.99,0.999$ give errors
$0.04,0.2,2,20$. At $r=1$ the channel is unidentifiable. Exact full rank near
$r=1$ does not establish useful precision.

Adding a third response with $r=2$ to the baseline and $r=1/2$ gives, at the
same three error radii $0.01$, optimal weights $(0,-2/3,2/3)$ and error $1/75$.
The dual $h=(1/75,-1/60)$ proves optimality. This improvement presumes that
the extra measurement does not worsen the existing error bounds; a fixed
physical budget may require a different comparison.

Unknown selectivity is a decisive control: $(\theta_1,\theta_2,r)=(3,1,1/2)$
and $(2,2,1/4)$ both give $(y_0,y_1)=(4,5/2)$. Inserting a nominal $r$ into a
known-matrix certificate does not validate it when $r$ is unknown.

\section{From one record to a finite-error experiment}

\subsection{Known laws and affinity}

For known finite record laws define their Bhattacharyya affinity
\cite{caBhattacharyya1943}
\[
 \rho(P_0,P_1)=\sum_z\sqrt{P_0(z)P_1(z)}.
\]
Cauchy--Schwarz gives $0\leq\rho\leq1$, with equality at one exactly when the
two laws agree. For $M$ independent records under the same fixed model,
the likelihood-ratio decision rule has equal-prior average error at most
\begin{equation}
 r_M^*\leq\tfrac12\rho(P_0,P_1)^M.
 \label{precision:affinity}
\end{equation}
To prove this, use $\min(p,q)\leq\sqrt{pq}$ in
Equation~\eqref{precision:tv-risk}; affinity multiplies over product laws.
An outward rational upper bound $\overline\rho<1$ therefore certifies a
sufficient record count. An upper enclosure at least one certifies no such
count at that precision. It need not mean that the actual laws are equal.

The product-law premise requires independent preparations or a justified
independent-record construction. If a nuisance parameter is fixed but unknown,
it stays fixed in every factor. Averaging its law independently at every record
models a different experiment. Full common support also prevents perfect
finite-record classification, even when the affinity bound becomes very small.

\subsection{A robust bounded-score witness}

Let $f(z)\in[0,1]$ be chosen before acquisition. It may be an event indicator.
Suppose all nominal models in target class $+$ satisfy
$\mathbb E f\geq a_0$, and all in class $-$ satisfy $\mathbb E f\leq b_0$.
Suppose actual one-record laws are within total variation $\eta_+$ or $\eta_-$
of their respective nominal laws. Since bounded-score expectations change by
at most total variation, define
\begin{equation}
 a=\max(0,a_0-\eta_+),\quad
 b=\min(1,b_0+\eta_-),\quad \gamma=a-b.
 \label{precision:robust-gap}
\end{equation}

\begin{proposition}[A uniform finite-record guarantee]
\label{precision:score-test}
If $\gamma>0$, classify as $+$ when the mean of $M$ scores exceeds
$(a+b)/2$, and as $-$ otherwise. For independent records under each fixed
actual model, each conditional error probability is at most
\[
 \exp(-M\gamma^2/2).
\]
In particular, $M\geq2\log(1/\alpha)/\gamma^2$ gives each error at most
$\alpha$. The same bound holds for a dependent sequence of scores in $[0,1]$
if every full-history conditional mean is at least $a$ under class $+$ or at
most $b$ under class $-$.
\end{proposition}
\begin{proof}
Under $+$, the rejection threshold is at least $\gamma/2$ below the mean.
Hoeffding's inequality gives the result, and the argument for $-$ is symmetric
\cite{caHoeffding1963}. For the conditional version, Hoeffding's bounded-variable
moment-generating-function lemma gives, under $+$,
\[
 \mathbb E[\exp(-t(X_i-a))\mid\mathcal F_{i-1}]
 \leq\exp(t^2/8),\qquad t\geq0.
\]
Iterate conditional expectations, apply Markov's inequality to the sum, and
minimize
\[
 \exp(-tM\gamma/2+Mt^2/8)
\]
at $t=2\gamma$.
\end{proof}

A nonpositive robust gap means that this certificate fails. It does not prove
that every possible protocol fails. A one-record separating score can also fail
for composite hypotheses even when repeated records identify them: a fixed
Bernoulli parameter in $\{0.2,0.8\}$ cannot be separated by a single bounded
score from parameter $0.5$, because the latter law is the average of the first
two. Repeated independent observations of the same fixed parameter separate
the classes through the sample mean.

The inverse-square bound is sufficient, not a universal sample-complexity law.
For Bernoulli probabilities $0$ and $\varepsilon$, deciding for the second
model after the first success has zero error under the first and error
$(1-\varepsilon)^M$ under the second. A fixed error then needs order
$\varepsilon^{-1}$ records. Interior alternatives such as
$1/2\pm\varepsilon$ behave differently.

\section{Calibration must cover the experiment that is executed}

For a classical hidden-state instrument, calibration concerns the joint
distribution of the next hidden state and observed outcome. Suppose the initial
state laws differ by at most $\eta_0$ in total variation and, at step $k$, every
admissible state, retained history and selected action has joint-kernel error
at most $\eta_k$. Require $0\leq\eta_k\leq1$, including $k=0$; any looser
total-variation bound can first be truncated at one. Under a common causal policy,
\begin{equation}
 \operatorname{TV}(P_{\mathrm{actual}}^\pi,P_{\mathrm{nominal}}^\pi)
 \leq 1-(1-\eta_0)\prod_{k=1}^H(1-\eta_k)
 \leq\eta_0+\sum_{k=1}^H\eta_k.
 \label{precision:kernel-calibration}
\end{equation}
Couple the initial states maximally and use the same policy randomization. As
long as the histories and hidden states agree, the next coupled state/outcome
agrees with conditional probability at least $1-\eta_k$. Multiplication gives
the first bound; a union bound gives the second. Stopping can be padded with an
identical absorbing action. Bounds on detector marginals alone are insufficient:
two instruments can give the same immediate output while placing the hidden
state in different conditions for the next experiment.

For quantum instruments, represent each entire outcome-retaining operation as a
completely positive trace-preserving map to a classical--quantum output. An
initial trace-distance bound $\eta_0$ and half-diamond-norm bounds $\eta_k$
give the additive analogue of \eqref{precision:kernel-calibration}. A telescoping
hybrid argument inserts one actual operation at a time; trace-distance
contractivity, including ancillas and later feedback channels, bounds its
contribution by $\eta_k$ \cite{caWatrous2018}. The classical multiplicative
coupling bound is not asserted for arbitrary quantum channels. A norm on an
unconstrained signed matrix representation is not by itself a physical channel
calibration certificate.

There are two distinct routes for propagating repeated-record errors. A uniform
one-record radius can be incorporated into the score gap before amplification.
Alternatively, an independently justified whole-dataset total-variation bound
$\Delta$ changes any nominal decision error by at most $\Delta$.
It is incorrect to add a one-record radius only once to a many-record decision
error. Under independent product-law models a radius $\eta$ yields the dataset
bound $1-(1-\eta)^M\leq M\eta$, which may be too loose to be useful.

\section{Choosing a protocol for the stated decision}

The shortest or lowest-charge exact witness minimizes a qualitative objective:
it obtains any nonzero distinction. It need not minimize the resources needed
for a prescribed error. A deeper single record may improve its contrast enough
to reduce total repetitions, or its extra control and reset costs may outweigh
that improvement. An experiment report should compare the full decision at
matched resource budgets.

For a finite candidate family $j=1,\ldots,K$, let $x_j$ be the unnormalized
hidden-state row vector whose mass $m_j=x_j\mathbf1$ includes the candidate's
prior and the probability of the observed history. Each candidate has a target
label $g_j$. If an action $a$ yields outcome $y$, the update is
$x_j'=x_jM_{j,a,y}$. A fixed nuisance value is represented within the index $j$;
it is never resampled between updates. Let
\[
 L(x)=\sum_jm_j-\max_v\sum_{j:g_j=v}m_j
\]
be the Bayes error mass incurred by stopping and choosing a target label. For
two single-model target classes this is $\min(m_0,m_1)$.

\begin{proposition}[Finite-horizon Bayes design]
\label{precision:bellman}
Fix finite classical instruments with nonnegative kernels satisfying
$\sum_y M_{j,a,y}\mathbf1=\mathbf1$, a common finite action/outcome library,
nonnegative integer charges $c(a)$, remaining budget $B$, and a finite remaining
action horizon $h$. A finite grammar state
$q$ determines legal actions, a transition $q'=\Phi(q,a)$, and whether stopping
is permitted. Assign every candidate a strictly positive prior, so zero-mass
branches are impossible under every candidate and may safely be ignored when
enforcing feasibility. The minimum Bayes error mass satisfies
\begin{equation}
 V_{h,B,q}(x)=\min\left\{
 \begin{array}{l}
 L(x)\quad\hbox{if stopping is permitted},\\[1mm]
 \displaystyle\min_{\substack{a\ \mathrm{legal},\ c(a)\leq B\\h>0}}
       \sum_y V_{h-1,B-c(a),\Phi(q,a)}
       ((x_jM_{j,a,y})_{j=1}^K)
 \end{array}\right\}.
 \label{precision:bellman-equation}
\end{equation}
An unavailable alternative has value $+\infty$; a zero-mass branch has value
zero. The initial vectors are the supplied prior-weighted initial states.
\end{proposition}
\begin{proof}
At horizon zero only permitted stopping is possible. At a positive horizon, a
causal policy first stops or selects an available action. Conditional histories
of distinct outcomes are disjoint, so their unnormalized error masses add.
By induction, each continuation has risk at least the corresponding displayed
value. Selecting minimizing continuations attains that value. Randomizing the
first action forms a convex combination and cannot lower the minimum.
\end{proof}

This is standard belief-state dynamic programming \cite{caBellman1957}.
Rational probabilities and charges allow exact finite arithmetic. Finite horizon
is essential even with zero-charge actions. Search can still grow exponentially;
an implementation reaching its declared state limit has not certified an
optimum. This result optimizes the given priors, model library, target loss,
grammar and budget. Robust minimax design, unknown nuisance priors, new physical
controls and different energy accounting are different optimization problems.
An unmodeled candidate is not repaired by solving this recurrence more exactly.

\subsection{Fisher information answers a different question}

For a regular scalar Bernoulli model, local Fisher information is
$I(\theta)=[p'(\theta)]^2/[p(\theta)(1-p(\theta))]$. Consider two equally costly
probes, $p_A(\theta)=\theta^2$ and $p_B(\theta)=(1+\theta/2)/2$, at the target
alternatives $\theta=\pm1/2$. Probe $A$ has Fisher information $16/3$ at either
point, while probe $B$ has $4/15$. Yet $A$ is exactly blind to the sign, whereas
$B$ distinguishes it. Probe $A$ resolves a local change of magnitude, so there
is no contradiction. The example explains why maximizing local information,
minimizing Bayes classification risk, and minimizing energy for target precision
must be stated as separate objectives.

The next chapter adds two requirements for use outside a closed model exercise:
the option to withhold a target when the candidate predictions fail, and a
physical resource ledger whose uncertainty composes with the inference guarantee.
